\chapter{术语与缩写}
\begin{longtable}{p{.18\linewidth}p{.72\linewidth}}
\toprule 名称 & 本报告中的含义\\\midrule\endhead
VLM & 视觉语言模型（Vision-Language Model），接收图像和文字，通常输出文字或位置描述。\\
VLA & 视觉语言动作模型（Vision-Language-Action Model），将视觉、语言和机器人状态映射为动作。\\
ACT & 动作分块变换器（Action Chunking with Transformers），一次预测后续一段动作。\\
LoRA & 低秩适配（Low-Rank Adaptation），通过训练较少的附加参数适配已有模型。\\
RGB-D & 彩色图像与深度图的组合。仿真真值深度和真实传感器深度分别说明。\\
BF16 & 脑浮点16位格式（Brain Floating Point 16），一种16位浮点表示。\\
NF4 & 4位正规化量化格式（NormalFloat 4），用于减少模型权重存储。\\
ROS2 & 机器人操作系统第二代，提供节点通信、坐标变换和工具接口。\\
Nav2 & ROS2导航框架，负责路径规划、局部避障和目标到达。\\
MoveIt & 机械臂运动规划框架，调用逆运动学、碰撞检查和轨迹执行接口。\\
MuJoCo & 用于机器人动力学和接触过程的物理模拟器。\\
LIBERO & 用于语言条件机器人操作学习与评测的公开任务基准。\\
可供性 & 场景中适合执行某类动作的位置或区域，例如可抓取的把手。\\
闭环 & 执行后重新读取观察与状态，再据此决定下一步。\\
P50 / P95 & 第50和第95百分位数。这里用于描述时延分布，需同时给出样本量。\\
交换空间 & 内存不足时使用的后备空间。加载阶段与稳定推理阶段分开记录。\\
影子记录 & 只记录真实状态供事后评分，不用它驱动模型决策或停止。\\
IoU & 交并比：预测区域与参考区域交集面积除以并集面积。\\
严格／邻域命中 & 操作点是否落入目标掩码；邻域命中允许距离掩码边界不超过4像素。\\
MAE / L1 & 平均绝对误差；需注明动作、像素、米或角度等具体单位。\\
FP16 / FP32 / INT8 & 16位浮点、32位浮点与8位整数表示；量化范围随实验注明。\\
KV Cache & 注意力键值缓存，复用已生成文字的中间状态以减少重复计算。\\
EOS & 序列结束标记；未输出结束标记或结构不完整可能是生成长度截断。\\
TF & 坐标变换系统，维护相机、机械臂与世界之间的坐标关系。\\
URDF & 统一机器人描述格式，描述连杆、关节、几何和惯量。\\
KDL / OMPL & 运动学与采样式运动规划工具，本项目通过MoveIt接口调用。\\
PD控制 & 比例与微分反馈控制，根据位置和速度误差产生控制量。\\
mocap驱动 & 通过模拟器的运动捕捉目标直接设定末端参考，条件比真实关节力矩控制理想。\\
JIT & 即时编译；首次调用包含编译开销，与预热后的推理分开测量。\\
CUDA / CUPTI & NVIDIA GPU计算平台及其性能采样接口。\\
RSS / cgroup & 进程驻留内存与Linux进程组资源限制；不与共享GPU内存简单相加。\\
JSON / JSONL & 结构化文本格式；JSONL逐行保存独立记录，便于中断后恢复。\\
SHA256 & 文件内容摘要，用于核对输入、权重和素材是否一致。\\
McNemar检验 & 用于配对二分类结果的统计检验；小样本下不能只凭成功数差异判断显著性。\\
Wilson区间 & 二项比例置信区间的一种估计方式，反映有限样本下的不确定性。\\
动作块 & 一次预测多步动作；执行其中若干步后重新观察和推理。\\
低秩合并 & 将适配器增量合入基座权重；须核对张量轴顺序和未合并参考的输出。\\
\bottomrule\end{longtable}
